% Cuestión VI. El método que no era LSB1, LSB4 ni LSBI: datos agregados al final de madrid.bmp.

\textbf{El método.} En \texttt{madrid.bmp} el mensaje se agregó \emph{al final del archivo}, después del último píxel. El archivo mide
\dato{an-madrid-bytes} bytes, pero la cabecera más los píxeles ocupan \dato{an-madrid-bfsize}: sobran \dato{an-madrid-agregados} bytes de texto en claro
(listado~\ref{lst:q6-madrid}). Los píxeles no ocultan nada.

\lstinputlisting[caption={Bytes agregados al final de \texttt{madrid.bmp}},label={lst:q6-madrid}]{generado/an-madrid-texto.txt}

\textbf{¿Es eficaz?} No:
\begin{itemize}
\item Se detecta comparando el tamaño del archivo con el de la cabecera, o mirando el final con \texttt{strings} o un editor hexadecimal.
\item No hay ninguna protección: el texto está en claro.
\item Se pierde al volver a guardar o convertir la imagen, porque los programas sólo conservan los píxeles.
\end{itemize}
Su única ventaja es que no altera los píxeles, por lo que no lo detectan los análisis estadísticos de los bits bajos y su capacidad es ilimitada. Sirve, como
lo usó la cátedra, para dejar una pista, no para ocultar un secreto.
